% Part: first-order-logic
% Chapter: axiomatic-deduction
% Section: proving-things-quant

\documentclass[../../../include/open-logic-section]{subfiles}

\begin{document}

\olfileid{fol}{axd}{prq}

\olsection{\usetoken{P}{derivation} مع المكمّمات}

\begin{ex}
لننشئ !!a{derivation} لــ$(\lforall[x][!A(x)] \land
\lforall[y][!B(y)]) \lif \lforall[x][(!A(x) \land !B(x))]$.
% تصحيح مصدر (0117-I2): لم يصرّح الأصل بجِدّة الثابت اللازمة لتطبيق قاعدة المكمّمات.
ونختار في البرهان ثابتًا مميَّزًا جديدًا~$a\in\mathcal E$ لا يرد
في~$!A(x)$ ولا في~$!B(y)$.

لاحظ أولًا أن
\begin{align*}
  (\lforall[x][!A(x)] \land \lforall[y][!B(y)]) & \lif \lforall[x][!A(x)]\\
  \intertext{حالة من \olref[prp]{ax:land1}، وأن}
  \lforall[x][!A(x)] & \lif !A(a) \\
  \intertext{حالة من \olref[qua]{ax:q1}. ولذلك نعرف، بحسب \olref[pro]{prop:chain}، أن}
  (\lforall[x][!A(x)] \land \lforall[y][!B(y)]) & \lif !A(a) \\
  \intertext{!!{derivable}. وبالمثل، بما أن}
  (\lforall[x][!A(x)] \land \lforall[y][!B(y)]) & \lif \lforall[y][!B(y)] \qquad\text{و}\\
    \lforall[y][!B(y)] & \lif !B(a)\\
    \intertext{حالتان من \olref[prp]{ax:land2} و\olref[qua]{ax:q1}، على الترتيب، فإن}
    (\lforall[x][!A(x)] \land \lforall[y][!B(y)]) & \lif !B(a)\\
    % تصحيح مصدر (0117-I1): لا تكفي بديهية الوصل وتطبيقا MP مباشرةً؛ لا بد من تركيب المقدّم ثم استعمال بديهية توزيع اللزوم.
    \intertext{قابلة للاشتقاق بحسب \olref[pro]{prop:chain}. ولنسمِّ المقدّم المشترك $!P$. نصل $!P \lif !A(a)$ بالحالة المناسبة من \olref[prp]{ax:land3} بواسطة \olref[pro]{prop:chain}، فنحصل على $!P \lif (!B(a) \lif (!A(a) \land !B(a)))$. ثم نأخذ الحالة المناسبة من \olref[prp]{ax:lif2} ونجري~\MP{} مرتين، إحداهما مع $!P \lif !B(a)$، فنحصل على أن}
    (\lforall[x][!A(x)] \land \lforall[y][!B(y)]) & \lif (!A(a) \land !B(a))\\
    \intertext{قابلة للاشتقاق. ويمكننا الآن تطبيق \QR{} لنحصل على}
(\lforall[x][!A(x)] \land \lforall[y][!B(y)]) & \lif \lforall[x][(!A(x) \land !B(x))].
\end{align*}
\end{ex}

\end{document}
